\documentclass[11pt]{article}
\usepackage[margin=1in]{geometry}
\usepackage{amsmath,amssymb,amsthm}
\usepackage{array}
\usepackage{parskip}
\usepackage{booktabs}

\newtheorem{theorem}{Theorem}
\newtheorem{lemma}[theorem]{Lemma}
\newtheorem{corollary}[theorem]{Corollary}
\newtheorem{proposition}[theorem]{Proposition}
\theoremstyle{definition}
\newtheorem{definition}[theorem]{Definition}
\theoremstyle{remark}
\newtheorem{remark}[theorem]{Remark}

\newcommand{\N}{\mathbb{N}}
\newcommand{\Z}{\mathbb{Z}}
\newcommand{\Q}{\mathbb{Q}}
\newcommand{\PP}{\mathbb{P}}
\newcommand{\OK}{\mathcal{O}}
\newcommand{\SZ}{\mathbb{Z}_S}

\title{Disproof of conjecture \texttt{00000001737}}
\author{earthking11}
\date{2026-09-14}

\begin{document}
\maketitle

\section*{Verdict: FALSE}

We disprove conjecture \texttt{00000001737} as stated. The refutation is
unconditional and elementary: for the two-element prime set $S=\{2,3\}$ the
projective line minus three points $\PP^1\setminus\{0,1,\infty\}$ has
\emph{21} $S$-integral points, not at most $12$. Since
$21 > 12$, the conjectured maximum is false. We additionally verify that
$12$ is not a value attained by any two-element prime set $S$ with
$p<q\le 23$, and that the true maximum there is $21$, attained at
$S=\{2,3\}$. The final phrase of the conjecture, ``with the extremum exhausted
by explicit hyperelliptic covers'', is a non-sequitur: $\PP^1$ minus three
points has genus $0$ and is not a hyperelliptic (nor any) curve of genus
$\ge 1$; no hyperelliptic cover is involved.

\section{The conjecture, quoted verbatim}

From \texttt{conjectures/00000001737.md}:

\begin{quote}
\textbf{English.} Conjecture: The maximal possible number of $S$-integral
points of $\PP^1\setminus\{0,1,\infty\}$ with $S$ of two primes is $12$; a
tightening of Evertse's general bound, with the extremum exhausted by explicit
hyperelliptic covers.
\end{quote}

The conjecture is filed in both English and Chinese; the Chinese original is
reproduced verbatim (in Unicode) in \texttt{README.md}, and the two versions are
identical in content. The PDF font does not embed CJK glyphs, so the Chinese is
kept in the markdown copy rather than duplicated here.

\section{The definition of $S$-integral point}

The statement is only meaningful under the standard definition, and it is worth
fixing carefully because the naive reading is degenerate.

Let $K=\Q$, let $S=\{p_1,p_2\}$ be a finite set of primes, and let
\[
\OK_S \;=\; \Z[1/p_1,1/p_2] \;=\; \Z[1/p_1 p_2]
\]
be the ring of $S$-integers. An element of $\OK_S^\times$ is an \emph{$S$-unit}:
a non-zero rational $u$ with $|u|=p_1^{a}p_2^{b}$ for some $a,b\in\Z$.

\begin{definition}\label{def:sint}
The scheme $\PP^1\setminus\{0,1,\infty\}$ over $\Z$ is
$\operatorname{Spec}\Z[x,1/x,1/(1-x)]$. An \emph{$S$-integral point} is a
$\Z$-morphism $\operatorname{Spec}\OK_S\to\PP^1\setminus\{0,1,\infty\}$, i.e.\
a rational $x$ such that
\[
x \in \OK_S^\times, \qquad 1-x \in \OK_S^\times .
\]
Equivalently: $x$, $1/x$ and $1/(1-x)$ have no prime outside $S$ in numerator
or denominator. Note that ``$1/x$ integral'' forces $v_p(x)=0$ for
$p\notin S$.
\end{definition}

The integrality is required at the \emph{whole} divisor $\{0,1,\infty\}$. This
is the only reading under which the problem is a finite count. The naive
reading ``$x\ne 0,1$ with $x\in\OK_S$'' would already make the set infinite:
$\Z[1/6]$ is infinite, and every $x\in\Z[1/6]$ with $x\ne 0,1$ would qualify. So
the naive reading cannot be what is meant, and we use Definition~\ref{def:sint}
throughout.

\begin{remark}[classical form]
Writing $y=1-x$, Definition~\ref{def:sint} says precisely that $(x,y)$ is a
solution in $S$-units of
\[
x+y=1 .
\]
This is the \emph{$S$-unit equation}. By Mahler's theorem (1933) it has
finitely many solutions, effective bounds going back to Gy\H{o}ry (1979); the
deep inputs (linear forms in logarithms, or lattice reduction) are needed only
for \emph{exact} computation. For the refutation we need only a lower bound,
which is elementary.
\end{remark}

\section{The decisive witness: $S=\{2,3\}$}

\begin{theorem}\label{thm:21}
For $S=\{2,3\}$ the set of $S$-integral points of
$\PP^1\setminus\{0,1,\infty\}$ has exactly $21$ elements, namely
\[
\begin{aligned}
&-8,\ -3,\ -2,\ -1,\ -\tfrac12,\ -\tfrac13,\ -\tfrac18,\\
&\tfrac19,\ \tfrac14,\ \tfrac13,\ \tfrac12,\ \tfrac23,\ \tfrac34,\
 \tfrac89,\ \tfrac98,\ \tfrac43,\ \tfrac32,\ 2,\ 3,\ 4,\ 9 .
\end{aligned}
\]
\end{theorem}

\begin{proof}[Verification]
Each listed $x$ has $x$ and $1-x$ of the form $\pm 2^{a}3^{b}$ with
$a,b\in\Z$; the corresponding pairs $(x,1-x)$ are
\[
\begin{array}{llll}
(-8,9) & (-3,4) & (-2,3) & (-1,2)\\
(-\tfrac12,\tfrac32) & (-\tfrac13,\tfrac43) & (-\tfrac18,\tfrac98) &
 (\tfrac19,\tfrac89)\\
(\tfrac14,\tfrac34) & (\tfrac13,\tfrac23) & (\tfrac12,\tfrac12) &
 (\tfrac23,\tfrac13)\\
(\tfrac34,\tfrac14) & (\tfrac89,\tfrac19) & (\tfrac98,-\tfrac18) &
 (\tfrac43,-\tfrac13)\\
(\tfrac32,-\tfrac12) & (2,-1) & (3,-2) & (4,-3)\\
\multicolumn{4}{l}{(9,-8)}
\end{array}
\]
Every entry of every pair is a unit of $\Z[1/6]$, so all $21$ rationals are
$S$-integral points by Definition~\ref{def:sint}. They are pairwise distinct,
so the count is at least $21$. The exact count $21$ (rather than merely
$\ge 21$) is confirmed by exhaustive enumeration: for every exponent bound
$B\ge 5$ the enumeration of $x=\pm 2^{a}3^{b}$ with $|a|,|b|\le B$ satisfying
Definition~\ref{def:sint} returns the same $21$ elements; the saturated list is
stable for all $B\in\{5,10,20,40,60,80,100,120,150\}$. See
\texttt{reproduce.py} and Section~\ref{sec:formal}.
\end{proof}

\begin{corollary}[the conjecture is false]\label{cor:false}
The maximal possible number of $S$-integral points of
$\PP^1\setminus\{0,1,\infty\}$ over two-element prime sets is at least $21$.
In particular it is not $12$, and the assertion ``the maximal possible number
$\dots$ is $12$'' is false.
\end{corollary}

\begin{proof}
Theorem~\ref{thm:21} gives $21$ points for the admissible two-element prime set
$S=\{2,3\}$, and $21>12$.
\end{proof}

\section{The maximum over all two-element prime sets}

To make the failure robust, we computed the exact count for every pair of primes
$p<q\le 23$ by the same saturated enumeration.

\begin{center}
\begin{tabular}{l r}
\toprule
$S$ & number of $S$-integral points \\
\midrule
$\{2,3\}$  & $21$ \\
$\{2,5\}$, $\{2,7\}$, $\{2,17\}$ & $9$ \\
$\{2,11\}$, $\{2,13\}$, $\{2,19\}$, $\{2,23\}$ & $3$ \\
any $\{p,q\}$ with $p,q$ both odd ($3\le p<q\le 23$) & $0$ \\
\bottomrule
\end{tabular}
\end{center}

The maximum over all these sets is $21$, attained only at $S=\{2,3\}$. No
two-element prime set in the range has exactly $12$ points: the observed counts
are $0,3,9,21$ only. The same is expected in general: $S=\{2,3\}$ is the
extremal two-element set for the $S$-unit equation.

\begin{remark}[cross-reference]
The values for $S=\{\text{first }n\text{ primes}\}$ are OEIS
\textbf{A362567}, ``Number of rational solutions to the $S$-unit equation
$x+y=1$, where $S=\{\text{prime}(i):1\le i\le n\}$'', whose data begin
\[
a(0)=0,\quad a(1)=3,\quad a(2)=21,\quad a(3)=99,\ \dots
\]
with offset $0$. Its a(2)${}=21$ is exactly our count, and its EXAMPLE section
lists precisely the $21$ solutions exhibited above. Our enumeration reproduces
A362567's a(2) independently and from scratch.
\end{remark}

\section{Comparison with Evertse's bound}

Evertse's general bound for the number of solutions of the $S$-unit equation is
\emph{exponential} in $|S|$; the commonly quoted form is
\[
N(S) \;\le\; 3\cdot 7^{\,2|S|+3}.
\]
For $|S|=2$ this is
\[
3\cdot 7^{7} = 3\cdot 823543 = 2\,470\,629,
\]
which is not close to $12$. It is also not close to $21$. So:

\begin{itemize}
\item $12$ is not Evertse's bound for $|S|=2$; the true Evertse-type bound is
  $2\,470\,629$.
\item $12$ is not Evertse's bound for any small $|S|$ either: for $|S|=1$ the
  bound is $3\cdot 7^{5}=50409$ (and the true value is $a(1)=3$), and for
  $|S|\ge 2$ the bound is already in the millions. In fact $12$ is not a
  standard bound in this circle of ideas at all.
\item Therefore ``a tightening of Evertse's general bound'' to $12$ cannot be
  obtained by any known bound: the claim is false on both counts, as a maximum
  (it is $21$ for $|S|=2$) and as an upper bound (Evertse's is exponential,
  not $12$).
\end{itemize}

The genuine exact values ($a(1)=3$, $a(2)=21$, $a(3)=99$) are far below
Evertse's bound; that gap is exactly why exact $S$-unit computations are hard,
and it is unrelated to the number $12$.

\section{The ``hyperelliptic covers'' clause}

The conjecture ends: ``with the extremum exhausted by explicit hyperelliptic
covers''. This is a non-sequitur for the objects at hand.

\begin{itemize}
\item $\PP^1\setminus\{0,1,\infty\}$ is a curve of genus $0$: the projective
  line has genus $0$ and removing points does not change the genus. Genus-$0$
  curves are \emph{not} hyperelliptic in the usual sense (no degree-$2$ map to
  $\PP^1$ of the required kind exists over $\Q$ with the relevant ramification
  — the only genus-$0$ case is the base itself). There is no hyperelliptic
  curve and no cover to speak of.
\item The $S$-unit equation $x+y=1$ is a problem on the torus
  $\mathbb{G}_m^2$, not on a hyperelliptic curve; its finiteness comes from
  Mahler's theorem and its effective resolution from linear forms in
  logarithms, not from hyperelliptic covering arguments.
\item The extremal set for $|S|=2$ is $S=\{2,3\}$, with $21$ points found by
  direct enumeration of $S$-units; no hyperelliptic cover is needed or used.
\end{itemize}

We therefore record the clause as mathematically unrelated to the claim, and
note that it does not rescue the conjecture.

\section{Framing: this is a lower-bound witness}\label{sec:framing}

We state the logical shape of the refutation precisely so as not to
overclaim.

\begin{itemize}
\item We prove \emph{existence} of $21$ pairwise-distinct $S$-integral points
  for $S=\{2,3\}$. Hence the maximum over two-element prime sets satisfies
  $\max \ge 21$.
\item Since $21 > 12$, the conjectured maximum of $12$ is refuted. This is all
  that is needed; an upper bound, or an exact maximum over all $S$, is not
  required.
\item We do \emph{not} claim to compute the exact maximum over all two-element
  prime sets. We do show that for $S=\{2,3\}$ the count is exactly $21$ by
  saturation of the enumeration, and for $p<q\le 23$ we list the exact counts,
  whose maximum is $21$. A proof that no larger set $S$ does better would
  require the $S$-unit theorem / linear forms in logarithms and is beyond this
  elementary witness.
\item Likewise we do not compute exact values for general $|S|$; those are the
  deep computations of A362567.
\end{itemize}

\section{Reproducibility}\label{sec:formal}

The exhaustive check is carried out by \texttt{reproduce.py} using the Python~3
standard library only (no \texttt{sympy}), with exact arithmetic via
\texttt{fractions.Fraction}:

\begin{quote}\small\ttfamily
\$ python3 reproduce.py\\
... \textit{(enumeration of $\pm 2^{a}3^{b}$ with exponent bound $60$;}\\
\textit{the full 21-element list; saturation across bounds $5,\dots,150$;}\\
\textit{counts for all prime pairs $p<q\le 23$; the checks)}\\
PASS: 21 distinct S-integral points exist for $S = \{2,3\}$;\\
\hspace*{1.5em}21 $>$ 12, so conjecture 00000001737 is FALSE.
\end{quote}

The refutation is also formalised in Lean~4 under \texttt{lean4/} (core Lean
only, \texttt{import Std}, no Mathlib, no \texttt{sorry}). The main statements
are
\begin{quote}\small\ttfamily
theorem isSmooth\_sound \{n : Nat\} (h : IsSmooth n = true) :\\
\hspace*{2em}\(\exists\) a b : Nat, n = 2 \^{} a * 3 \^{} b\\
theorem sIntegral\_sound \{p : Int\} \{q : Nat\} (h : SIntegral p q) :\\
\hspace*{2em}0 < q \(\wedge\) p \(\ne\) 0 \(\wedge\) (\(\exists\) a b, p.natAbs = 2\^{}a*3\^{}b)\\
\hspace*{2em}\(\wedge\) (\(\exists\) a b, q = 2\^{}a*3\^{}b) \(\wedge\) p \(\ne\) q \(\wedge\)\\
\hspace*{2em}(\(\exists\) a b, ((q:Int)-p).natAbs = 2\^{}a*3\^{}b)\\
theorem conjecture\_00000001737\_false :\\
\hspace*{2em}pts.length = 21 \(\wedge\) (\(\forall\) x \(\in\) pts, SIntegral x.1 x.2) \(\wedge\)\\
\hspace*{2em}(\(\forall\) x \(\in\) pts, \(\forall\) y \(\in\) pts, x \(\ne\) y \(\to\)\\
\hspace*{4em}x.1*(y.2:Int) \(\ne\) y.1*(x.2:Int)) \(\wedge\)\\
\hspace*{2em}(12 : Nat) < pts.length
\end{quote}
The rationals are encoded as pairs \texttt{(p : Int, q : Nat)} (lowest terms)
because \texttt{decide} cannot reduce \texttt{Rat} (its normalisation is
opaque); comparisons use cross-multiplication. \texttt{sIntegral\_sound} is the
faithfulness bridge from the executable Boolean test to the mathematical
definition. See \texttt{lean4/README.md} for the statement table, the proof
strategy, and the axiom audit (which reports only \texttt{propext} and
\texttt{Quot.sound}, and no \texttt{sorryAx}).

\section*{Status against the submission rules}

Rule 3 requires a LaTeX source, a PDF document, and a Lean~4 project.

\begin{itemize}
\item \textbf{LaTeX source} --- \texttt{main.tex}, a standalone \texttt{article}
  using \texttt{geometry}, \texttt{amsmath}, \texttt{amssymb}, \texttt{amsthm},
  \texttt{array}, \texttt{parskip}, \texttt{booktabs}; compiles with
  \texttt{tectonic main.tex}.
\item \textbf{PDF} --- at \texttt{build/main.pdf}, produced by
  \texttt{tectonic --outdir build main.tex}.
\item \textbf{Lean~4 project} --- \texttt{lean4/} with \texttt{lean-toolchain}
  pinned to \texttt{leanprover/lean4:v4.33.1}, \texttt{lakefile.toml}
  (library \texttt{Main}, no dependencies), \texttt{Main.lean},
  \texttt{Check.lean} (\texttt{\#print axioms}), and \texttt{lean4/README.md}.
  \texttt{lake build} exits $0$ and the audit reports no \texttt{sorryAx}.
\end{itemize}

\end{document}
